\chapter{Reflexionsbericht zur Nutzung von KI-Tools}
\label{chap:reflexionsbericht}

Gemäß den Betreuungsrichtlinien von Prof.\,Dr.\,Simon Kloker sowie \textcite{westermann2026ki} reflektiert dieser Bericht den Einsatz generativer KI-Werkzeuge im Forschungsprozess. Die wissenschaftliche Gesamtverantwortung, Konzeption, Validierung und finale Textfassung lagen durchgehend beim Autor.

\subsubsection*{1. Kriteriengeleitete Werkzeugauswahl}
Für Software-Prototyping, Datenexploration der Formcycle-Exporte und sprachliche Überarbeitung wurde \textit{Google Gemini} in einer agentischen Entwicklungsumgebung genutzt (großes Kontextfenster für \texttt{.fcp}/\texttt{.fcexport}-Modelle, lokaler Skriptaufruf, Schema-Treue). Zur datenschutzkonformen Evaluation der semantischen Prüfstufe (Art.\,44\,ff. DSGVO) wurden fünf Open-Weight-Modelle auf der HPC-Infrastruktur der \textit{Academic Cloud} (GWDG) verglichen (\textit{Llama 3.1 8B}, \textit{Qwen 3.8 27B}, \textit{Qwen 3 30B A3B}, \textit{GLM 5.3 Flash}, \textit{DeepSeek V4 Flash}).

\subsubsection*{2. Häufigkeit, Nutzen und Grenzen}
Hohen Nutzen stiftete die KI beim iterativen Engineering der Python-Prüfpipeline (flache \texttt{.fcp}-\texttt{items}-Auflösung über \texttt{parentid}, rekursive Workflow-Traversierung, Regex-Katalog) sowie beim UI-Prototyping des Dashboards. Das Architekturdiagramm (Abbildung~\ref{fig:pipeline_datenfluss}) wurde vom Autor eigenständig in \textit{draw.io} erstellt; lediglich einzelne Knotenbegriffe wurden KI-gestützt geschärft. Ungeeignet waren Standard-Chat-Prompts auf unkomprimierten Rohdaten: Ohne die Token-Kompression ($>90\,\%$) traten Kontextverluste (\textit{Lost in the Middle}) und halluzinierte Feldattribute auf.

\subsubsection*{3. Technische und wissenschaftliche Herausforderungen}
\enlargethispage{4\baselineskip}
Drei Herausforderungen prägten den Einsatz: Erstens fungierte die KI bei der Scoring-Formel (Gl.~\ref{eq:compliance_score}) als Diskussionspartner für Randwertstabilität, während die Schwellenwerte theoriegeleitet nach \textcite{nielsen1994severity} festgelegt wurden. Zweitens offenbarte \textit{Llama 3.1 8B} Zählfehler bei negativen Längenvorgaben (\texttt{HWG-LANG-02}), weshalb quantitative Vor-Metriken deterministisch in Python berechnet werden. Drittens erzeugten Sprachmodelle in frühen Tests erfundene Literaturhinweise; daher wurden alle 18 Quellen eigenständig in Fachdatenbanken (ACM, IEEE, ScienceDirect, SpringerLink) recherchiert und über DOIs verifiziert.

\subsubsection*{4. Mehrstufige Validierung (Vier-Augen-Prinzip)}
Alle KI-gestützten Artefakte wurden nach \textcite{recker2013scientific} mehrstufig geprüft: (1)~Regressionstests der Python-Engine an allen 34 \texttt{.xfp}-Archiven inklusive Fail-Safe-Prüfung bei fehlender API-Verbindung, (2)~manuelles JSON-Audit der synthetischen Testformulare (\texttt{synthetic\_compliant\_form.json} sowie 19 gezielt injizierte Regelverstöße in \texttt{synthetic\_\allowbreak noncompliant\_\allowbreak form.json}) und (3)~Abgleich aller Klickpfade und Regeln gegen BITV 2.0, WCAG 2.1, DIN EN ISO 9241-110 und Formcycle v8.

\subsubsection*{5. Erforderliche Fach- und Methodenkompetenzen}
Produktives Context-Engineering setzte fundierte Fachkenntnisse in Barrierefreiheit, Ergonomie und Datenmodellierung voraus, um disjunkte Fehlerklassen zu spezifizieren, Fehlklassifikationen (etwa bei \texttt{HWG-FEEDBACK-01}) zu erkennen und das Scoring deterministisch vom probabilistischen Sprachmodell zu entkoppeln.

\subsubsection*{6. Einfluss auf Denk- und Schreibprozess}
Die Entlastung von Boilerplate-Codierung schuf Freiraum für die konzeptionelle DSR-Arbeit und empirische Korpusanalyse. Dem Risiko eines \textit{Automation Bias} wurde durch kritische Gegenprüfung aller Modellausgaben gegen die realen Rohdaten und eine eigenständige Argumentationsführung begegnet.

\subsubsection*{7. Leitlinien für künftige Arbeiten}
(1)~\textit{Aufgabentrennung:} Syntax-, Struktur- und Rechenoperationen gehören in deterministischen Code, Sprachmodelle in die semantische Bewertung. (2)~\textit{Context-Engineering vor Modellgröße:} Komprimierte Extrakte mit Negativ-Constraints ermöglichten kompakten Open-Weight-Modellen ($\ge 27$B) am Testformular fehlerfreie Ergebnisse. (3)~\textit{Souveräne Hochschul-KI:} Akademische Cloud-Infrastrukturen ermöglichen leistungsfähige KI-Governance ohne Preisgabe personenbezogener Verwaltungsdaten.
